% Part: first-order-logic
% Chapter: sequent-calculus
% Section: derivations

\documentclass[../../../include/open-logic-section]{subfiles}

\begin{document}

\iftag{FOL}
      {\olfileid{fol}{seq}{der}}
      {\olfileid{pl}{seq}{der}}

\olsection{\usetoken{P}{derivation}}

\begin{explain}
ప్రారంభ సీక్వెంట్ ఎలా ఉంటుందో చెప్పాం; నిగమన నియమాలనూ ఇచ్చాం.
వీటి నుంచి సీక్వెంట్ కలనంలోని \tetoken{వ్యుత్పత్తులు}{!!^{derivation}s} ఆగమనాత్మకంగా
ఉత్పన్నమవుతాయి: ప్రతి \tetoken{వ్యుత్పత్తి}{!!{derivation}} తనంతట తానే ఒక ప్రారంభ సీక్వెంట్
అయినా అవుతుంది, లేదా ఒకటి లేదా రెండు \tetoken{వ్యుత్పత్తులు}{!!{derivation}s} తర్వాత ఒక
నిగమనం ఉండే రూపంలో ఉంటుంది.
\end{explain}

\begin{defn}[$\Log{LK}$ \tetoken{వ్యుత్పత్తి}{!!{derivation}}]
ఒక సీక్వెంట్~$S$కు \emph{$\Log{LK}$-\tetoken{వ్యుత్పత్తి}{!!{derivation}}} అంటే కింది
షరతులను తీర్చే సీక్వెంట్ల పరిమిత వృక్షం:
\begin{enumerate}
\item వృక్షంలో అన్నిటికంటే పైనున్న సీక్వెంట్లు ప్రారంభ సీక్వెంట్లు.
\item వృక్షంలో అన్నిటికంటే కిందనున్న సీక్వెంట్~$S$.
\item $S$ తప్ప వృక్షంలోని ప్రతి సీక్వెంట్, తనకు నేరుగా కింద నిలిచే
  నిష్కర్ష గల నిగమన నియమాన్ని సరిగ్గా వర్తింపజేయడంలో ఒక పూర్వాధారం.
\end{enumerate}
అప్పుడు $S$ను ఆ \tetoken{వ్యుత్పత్తి}{!!{derivation}} యొక్క \emph{అంత్య సీక్వెంట్} అనీ,
$S$ను \emph{$\Log{LK}$లో \tetoken{వ్యుత్పాదించదగిన}{!!{derivable}}} (లేదా
$\Log{LK}$-\tetoken{వ్యుత్పాదించదగిన}{!!{derivable}}) అనీ అంటాం.
\end{defn}

\begin{ex}
ప్రతి ప్రారంభ సీక్వెంట్, ఉదాహరణకు $!C \Sequent !C$, ఒక
\tetoken{వ్యుత్పత్తి}{!!a{derivation}}. దీనికి, ఉదాహరణకు, $\LeftR{\Weakening}$ నియమాన్ని
వర్తింపజేసి కొత్త \tetoken{వ్యుత్పత్తిని}{!!{derivation}} పొందవచ్చు:
\begin{prooftree}
\Axiom$ \Gamma \fCenter \Delta $
\RightLabel{\LeftR{\Weakening}}
\UnaryInf$ !A, \Gamma \fCenter \Delta$
\end{prooftree}
అయితే ఈ నియమం సాధారణంగా వర్తించేది: నియమంలోని $!A$ స్థానంలో ఏ
\tetoken{వాక్యమైనా}{!!{sentence}}, ఉదాహరణకు $!D$ను కూడా పెట్టవచ్చు. పూర్వాధారం మన
ప్రారంభ సీక్వెంట్ $!C \Sequent !C$తో సరిపోలితే, $\Gamma$, $\Delta$
రెండూ కేవలం $!C$ అని అర్థం; అప్పుడు నిష్కర్ష $!D, !C \Sequent !C$
అవుతుంది. కాబట్టి కిందిది ఒక \tetoken{వ్యుత్పత్తి}{!!a{derivation}}:
\begin{prooftree}
\Axiom$ !C \fCenter !C $
\RightLabel{\LeftR{\Weakening}}
\UnaryInf$ !D, !C \fCenter !C$
\end{prooftree}
ఇప్పుడు ఎడమవైపు ఉన్న రెండు \tetoken{వాక్యాలు}{!!{sentence}s} స్థానాలను మార్చనిచ్చే మరొక
నియమాన్ని, ఉదాహరణకు $\LeftR{\Exchange}$ను, వర్తింపజేయవచ్చు. అందువల్ల
కిందిది కూడా సరైన \tetoken{వ్యుత్పత్తి}{!!{derivation}}:
\begin{prooftree}
\Axiom$ !C \fCenter !C $
\RightLabel{\LeftR{\Weakening}}
\UnaryInf$ !D, !C \fCenter !C$
\RightLabel{\LeftR{\Exchange}}
\UnaryInf$ !C, !D \fCenter !C$
\end{prooftree}
ఈ నియమ వర్తనలో, ఇచ్చిన నియమ రూపం
\begin{prooftree}
\Axiom$ \Gamma, !A, !B, \Pi \fCenter \Delta $
\RightLabel{\LeftR{\Exchange}}
\UnaryInf$ \Gamma, !B, !A, \Pi \fCenter \Delta,$
\end{prooftree}
అయితే, $\Gamma$, $\Pi$ రెండూ శూన్య క్రమాలు, $\Delta$ అనేది $!C$;
$!A$, $!B$ పాత్రలను వరుసగా $!D$, $!C$ పోషిస్తాయి. దాదాపు ఇదే
విధంగా కిందిది కూడా ఒక \tetoken{వ్యుత్పత్తి}{!!a{derivation}} అని చూస్తాం:
\begin{prooftree}
\Axiom$ !D \fCenter !D $
\RightLabel{\LeftR{\Weakening}}
\UnaryInf$ !C, !D \fCenter !D$
\end{prooftree}
ఇప్పుడు ఈ రెండు \tetoken{వ్యుత్పత్తులను}{!!{derivation}s} తీసుకొని $\RightR{\land}$తో
కలపవచ్చు. ఆ నియమం:
\begin{prooftree}
\Axiom$\Gamma \fCenter \Delta, !A$
\Axiom$ \Gamma \fCenter \Delta, !B$
\RightLabel{\RightR{\land}}
\BinaryInf$ \Gamma \fCenter \Delta, !A \land !B $
\end{prooftree}
మన సందర్భంలో పూర్వాధారాలు, వాటితో ముగిసే \tetoken{వ్యుత్పత్తులలోని}{!!{derivation}s} చివరి
సీక్వెంట్లతో సరిపోలాలి. అందువల్ల $\Gamma$ అనేది $!C, !D$;
$\Delta$ శూన్యం; $!A$ అనేది $!C$; $!B$ అనేది $!D$. కాబట్టి నిగమనం
సరైనదైతే నిష్కర్ష $!C, !D \Sequent !C \land !D$ అవుతుంది.
\begin{prooftree}
\Axiom$ !C \fCenter !C $
\RightLabel{\LeftR{\Weakening}}
\UnaryInf$ !D, !C \fCenter !C$
\RightLabel{\LeftR{\Exchange}}
\UnaryInf$ !C, !D \fCenter !C$
\Axiom$ !D \fCenter !D $
\RightLabel{\LeftR{\Weakening}}
\UnaryInf$ !C, !D \fCenter !D$
\RightLabel{\RightR{\land}}
\BinaryInf$ !C, !D \fCenter !C \land !D $
\end{prooftree}
అవశ్యంగా, పూర్వాధారాల క్రమాన్ని తిప్పవచ్చు; అప్పుడు $!A$ అనేది
$!D$, $!B$ అనేది $!C$ అవుతాయి.
\begin{prooftree}
\Axiom$ !D \fCenter !D $
\RightLabel{\LeftR{\Weakening}}
\UnaryInf$ !C, !D \fCenter !D$
\Axiom$ !C \fCenter !C $
\RightLabel{\LeftR{\Weakening}}
\UnaryInf$ !D, !C \fCenter !C$
\RightLabel{\LeftR{\Exchange}}
\UnaryInf$ !C, !D \fCenter !C$
\RightLabel{\RightR{\land}}
\BinaryInf$ !C, !D \fCenter !D \land !C $
\end{prooftree}
\end{ex}

\end{document}
